\chapter{Model Architecture}

\section{Prior confined-environment ABMs}
CTB is designed to model a regulated, instrumented confined environment rather than an unobserved epidemic curve. Prior confined-environment models provide parts of this capability, but they
usually emphasize one layer at a time: individual contacts, spatial movement, route attribution, airflow, or a small intervention set. The distinction matters for cruise ships because the
empirical endpoint is already altered by case recognition, isolation, cleaning, food-service changes, and regulatory reporting. CTB therefore treats infection generation, observation, and
response as one coupled system.

Table~\ref{tab:confined-abm-comparison} summarizes the comparison used to position the model. The table includes agent-based models and closely related mechanistic ship or indoor-air models
because cruise-ship transmission lies at the intersection of individual mobility and multizone environmental transport. Kerkmann 2024 and Kou 2021 were not cited because the literature audit
could not verify those items as requested. The defensible novelty claim is narrow: among the accessible comparators reviewed here, none was found to combine individual route-weighted dose
response, CONTAM-style multizone HVAC coupling, multiple pathogens, VSP-style escalation, wearable biosurveillance, and operational decision logic in one cruise-ship framework.

{\small\setlength{\tabcolsep}{3pt}\begin{longtable}{p{0.14\textwidth}p{0.15\textwidth}p{0.22\textwidth}p{0.23\textwidth}p{0.16\textwidth}}
\caption{Prior confined-environment transmission models and their relationship to CTB.}\label{tab:confined-abm-comparison}\\
\toprule
Study & Environment and type & Transmission representation & Calibration or validation & Limitation relative to CTB \\
\midrule
\endfirsthead
\toprule
Study & Environment and type & Transmission representation & Calibration or validation & Limitation relative to CTB \\
\midrule
\endhead
\citet{Hunter2018OpenDataABM} & Open-data ABM for Irish towns & NetLogo spatial SEIR with co-location contacts, census and land-use driven agents, and stochastic disease periods & Tested against the 2012 Schull measles outbreak and transferred to multiple towns & No ship geometry, HVAC transport, route-specific dose accumulation, or regulatory response layer \\
\citet{Pan2026HealthyCruiseShip} & Reported cruise-ship ABM for healthy ship design & Spatial and operational interventions are indicated by metadata, but detailed mechanics were not available in the reviewed text & Not recoverable from accessible full text in the literature review & Potentially adjacent to CTB, but feature-level comparison could not be verified \\
\citet{Rankin2022DiamondPrincessABM} & NetLogo3D Diamond Princess ABM & Stochastic SEIR agents with passenger and crew movement, proximity contacts, and quarantine-restricted movement & Compared with the observed Diamond Princess attack fraction and a previous SEIR estimate & No explicit aerosol concentration, fomite route, dose-response, or dynamic instrument-driven response \\
\citet{Azimi2021DiamondPrincess} & Mechanistic Diamond Princess pathway model, not a conventional ABM & Markov-chain environmental states separating short-range, long-range, fomite, aerosol, and large-droplet pathways & Large scenario sweep accepted by epidemic-curve fit criteria & Aggregated population states rather than autonomous ship agents and operational protocols \\
\citet{Xia2023CONTAMCruiseShip} & CONTAM plus Wells--Riley risk on a passenger ship & Multizone airflow and airborne-virus risk under prescribed occupant assumptions & Applied to an operating-vessel case without outbreak-curve validation in the retrieved evidence & Strong air-transport layer, but not an individual outbreak model with clinical observation and SOP response \\
\citet{DOrazio2020RestartPublicBuildings} & Public-building ABM calibrated using Diamond Princess & Occupant movement, close-contact duration, masks, density reduction, and infectious-entry restrictions & Transmission coefficient tuned to the Diamond Princess epidemic segment and used for building scenarios & No ship operating model, multizone HVAC solver, or multi-pathogen dose framework \\
\citet{Zheng2016CruiseShipInterventions} & Cruise-ship individual contact network and SEIR-like model & Cabins, restaurants, shared spaces, air-change rate, HEPA, UVGI, masks, and quarantine & Validated against a previous cruise-ship influenza outbreak & Ship-specific intervention comparison without the surveillance cascade, multiple pathogens, or CTB's route-weighted dose ledger \\
\bottomrule
\end{longtable}}

This comparison motivates the architectural choice. The model must represent where people are, what route-specific dose they receive, which instrument would observe the event, and what
response would be triggered. A pure contact ABM cannot answer airflow questions. A pure CONTAM or Wells--Riley calculation cannot answer quarantine-compliance or diagnostic-delay questions.
A static intervention sweep cannot represent an escalation system in which the intervention exists only after the ship generates sufficient evidence.

\section{CTB architecture overview}
CTB is implemented in the public repository \url{https://github.com/bckirkup/Crusher_to_the_Bridge} and documented in code-grounded DeepWiki pages for the architecture overview, transmission
core, and reactive protocol engine.\footnote{DeepWiki pages used here are the indexed \url{https://deepwiki.com/bckirkup/Crusher_to_the_Bridge/1-architecture-overview},
\url{https://deepwiki.com/bckirkup/Crusher_to_the_Bridge/1.2-four-pathway-transmission-core}, and \url{https://deepwiki.com/bckirkup/Crusher_to_the_Bridge/1.6-reactive-protocol-engine}.} The
codebase uses a decoupled architecture: a central orchestrator owns the epoch loop and passes state among engines for mobility, infection biology, transport, observation, protocols, cost,
and telemetry. The design is intentionally modular because shipboard outbreaks couple processes with different natural units: individuals, zones, airflows, samples, SOPs, and resource
ledgers.

Each simulation is a discrete-time voyage. The present implementation uses hourly epochs, matching the transport bridge constant \texttt{HOURS\_PER\_EPOCH = 1.0} and the per-epoch
orchestrator design. Let $t=0,1,\ldots,T$ index epochs. At each epoch, CTB advances the following state vector:
\begin{equation}
\mathcal{S}(t)=\{X_i(t), L_i(t), B_i(t), C_i(t), M_{z,p}(t), F_{z,p}(t), A_{z,p}(t), R(t)\},
\end{equation}
where $X_i$ is the disease state of agent $i$, $L_i$ is location, $B_i$ is behavioral or compliance state, $C_i$ is cabin or class assignment, $M_{z,p}$ is airborne pathogen mass for
pathogen $p$ in zone $z$, $F_{z,p}$ is surface or food-associated mass, $A_{z,p}$ is the set of instrument observations available for zone or agent targets, and $R(t)$ is the regulatory or
SOP response state.

The epoch loop is ordered to avoid circular interpretation of a response. First, the Korkin-derived ship engine updates agent movement and disease progression. Second, pathogen shedding and
route-specific exposure are calculated from current locations and persistent environmental pools. Third, aerosol mass is transported natively or through the CONTAM branch. Fourth,
observation instruments sample agents, surfaces, air, wastewater, wearable streams, or clinical specimens. Fifth, stoplight logic converts raw instrument results into GREEN, AMBER, or RED
signals. Sixth, the protocol engine evaluates SOP triggers, applies modifiers to the next epoch, and debits activation or recurring costs. Finally, telemetry is serialized as simulation
history, artificial-lab-notebook records, contact tracing matrices, and resource ledgers.

Agents are initialized by \texttt{orchestrator\_init.py} from the platform spatial layout, role fractions, pathogen profiles, and voyage configuration. The initialization layer assigns
passenger or crew roles, home zones, cabin-mate sets, graywater collection zones, isolation-unit capacity, and pathogen seeding. Cabin grouping is explicit: agents assigned to a
\texttt{Cabin\_Corridor} are grouped into stateroom-size cohorts, and those cabin-mate IDs are later used to distinguish household-like exposure from non-mate corridor exposure during
confinement. This is a material distinction for ship simulations because quarantine does not remove a person from the ship air system and does not eliminate cabin-mate exposure.

The spatial model is a graph of zones rather than a continuous geometry. Zones carry attributes such as type, traffic class, deck, volume, cabin ventilation type, cabin size, maximum
occupancy, dining service type, and food-contamination multiplier. This representation is detailed enough to support cabin corridors, dining venues, galleys, medical spaces, public rooms,
engineering spaces, graywater collection, and HVAC nodes while still allowing large campaign sweeps. Campaign scale for the monograph is 69,750 simulation runs across 12 campaigns, 4
pathogens, and 4 vessel platforms, with 285 CDC VSP empirical outbreak records used as the calibration reference set.

\section{Agent state and pathogen seeding}
For pathogen $p$, agent $i$ has a disease state
\begin{equation}
X_{i,p}(t)\in\{S,E,I_{a},I_{s},R,Q\},
\end{equation}
where $S$ is susceptible, $E$ exposed, $I_a$ asymptomatic infectious, $I_s$ symptomatic infectious, $R$ recovered or no longer infectious, and $Q$ confined, quarantined, or isolated. The
implementation stores these states through the infection-dynamics bridge and pathogen-aware agent methods rather than a single global SIR label. Multiple concurrent pathogens are therefore
represented as pathogen-specific infection histories.

Pathogen seeding can occur at initialization or during the voyage. The initialization layer loads active pathogen profiles and infects configured initial agents or introduces pathogen mass
into environmental reservoirs. Mid-cruise introductions are handled as epoch events so that port visits, embarkation layers, or late infectious arrivals can be represented without changing
the core transmission equations. Let $I_{p}^{0}$ be the set of initially infected agents for pathogen $p$ and $J_p(t)$ be any introduction process during the voyage. Then the infected set
evolves by
\begin{equation}
\mathcal{I}_{p}(t+1)=\mathcal{I}_{p}(t)\cup \mathcal{N}_{p}(t)\cup J_p(t),
\end{equation}
where $\mathcal{N}_{p}(t)$ is the set of new infections produced by the dose-response step. This separation keeps imported cases distinct from secondary transmission, which is necessary for
estimating offspring distributions and fade-out probability.

Shedding is profile-specific. The common implementation form is a log-scale shedding curve adjusted by a fitted dose parameter:
\begin{equation}
V_{i,p}(t)=10^{s_p(\tau_{i,p}(t))-\doseadj_p},
\end{equation}
where $s_p(\tau)$ is the pathogen-specific log shedding trajectory at time since infection or symptom onset $\tau$, and $\doseadj_p$ maps biological shedding units to effective simulation
dose. This equation is the interface between pathogenesis and transmission. It allows norovirus, SARS-CoV-2, measles, Legionella, or other configured pathogens to share the same route ledger
while retaining different shedding timing, dose-response parameters, route weights, and environmental persistence.

\section{Four-pathway transmission core}
The transmission engine is called \texttt{TransmissionCore} in \texttt{engines/}\allowbreak\texttt{transmission\_core.py}. It began as a four-pathway engine with direct contact, short-range droplet, long-range
airborne transport, and fomite deposition. The current implementation retains those four primary pathways and extends them with food contamination and environmental-source pathways for
enteric and colonized-environment pathogens. For formal notation, define
\begin{equation}
\mathcal{R}=\{\mathrm{contact},\mathrm{droplet},\mathrm{airborne},\mathrm{fomite},\mathrm{food},
\mathrm{environmental}\}.
\end{equation}
When a pathogen profile does not enable food or environmental-source transmission, the corresponding dose terms are zero.

The route-weighted dose received by agent $i$ for pathogen $p$ in epoch $t$ is
\begin{equation}
D_{i,p}^{\mathrm{total}}(t)=\sum_{r\in\{\mathrm{contact},\mathrm{droplet},\mathrm{airborne},\mathrm{fomite},\mathrm{food}\}}
w_{p,r}\,D_{i,p,r}(t),
\label{eq:route-weighted-dose}
\end{equation}
with an additional environmental term when enabled. Here $w_{p,r}$ is the pathogen-specific transmission-route weight and $D_{i,p,r}$ is the raw route dose delivered by the pathway module.
The code resolves absent route weights to identity values, so a profile can opt into route weighting without changing the default behavior. The per-agent dose accumulator then applies
susceptibility multipliers before infection sampling.

The default infection probability uses the beta-Poisson form specified in the pathogen profile:
\begin{equation}
P_{i,p}(\mathrm{infection}\mid D)=1-\left(1+\frac{D_{i,p}^{\mathrm{total}}}{\beta_p}\right)^{-\alpha_p}.
\label{eq:beta-poisson}
\end{equation}
An exponential dose-response option is also implemented for pathogens or calibration settings that use a single rate constant:
\begin{equation}
P_{i,p}(\mathrm{infection}\mid D)=1-\exp(-\kappa_p D_{i,p}^{\mathrm{total}}).
\end{equation}
Only susceptible, non-immune agents are sampled. If infection occurs, the event is recorded with the dominant pathway, total dose, pathogen ID, zone, and pathway breakdown.

Direct contact is calculated inside each occupied zone. In density-dependent mode, the effective number of contacts for target agent $i$ in a zone with $N_z(t)$ occupants is sampled around
\begin{equation}
\bar{C}_{i,z}(t)=c_0\left(\frac{N_z(t)}{N_{\mathrm{ref}}}\right)^{\alpha_c}m_i(t),
\end{equation}
where $c_0$ is the base contact rate, $N_{\mathrm{ref}}$ is reference occupancy, $\alpha_c$ is the contact-density exponent, and $m_i(t)$ includes voyage contact scaling and crew service
multipliers in dining or galley zones. A hard maximum caps the draw. In legacy mode, contacts are drawn from the original average-contact pool. In heterogeneous-zone-dose mode, a mean-one
lognormal exposure factor
\begin{equation}
H_z\sim\exp\{\mathcal{N}(-\sigma_z^2/2,\sigma_z^2)\}
\end{equation}
multiplies within-zone dose so that expected mean exposure is preserved while crowded or service venues can generate stronger tails.

Short-range droplet exposure converts a fraction of shedding to same-room aerosol mass. The implementation uses a droplet aerosolization fraction and the clock-scaled inhaled air volume
$B_e=\operatorname{amount\_per\_epoch}(14.4\ {\rm m^3/day})$. For target $i$ in zone $z$,
\begin{equation}
D_{i,p,\mathrm{droplet}}(t)=\phi_{\mathrm{drop}}\,
\frac{\sum_{j\in\mathcal{I}_{p,z}(t)}V_{j,p}(t)}{V_z}\,
B_e\,q_i(t)\,b_z,
\end{equation}
where $\phi_{\mathrm{drop}}$ is the droplet aerosol fraction, $V_z$ is zone volume, $q_i(t)$ is the confinement factor, and $b_z$ is the balcony or zone ventilation
factor. The same shedding mass is also added to the zone aerosol pool for later transport.

Long-range airborne exposure uses the mass present in a target zone after native or CONTAM transport. If $M_{z,p}(t)$ is aerosol mass and $V_z$ is zone volume, then
\begin{equation}
D_{i,p,\mathrm{airborne}}(t)=\frac{M_{z,p}(t)}{V_z}\,B_e\,a_{\mathrm{HVAC}}(t)b_z,
\end{equation}
where $a_{\mathrm{HVAC}}$ is the active HVAC-airborne scalar injected by protocols. The explicit inhaled-volume term preserves the concentration-per-volume dependence, so larger zones dilute a fixed aerosol mass rather than receiving the same whole-zone dose.

Fomite exposure is persistent across epochs. Surface pools receive deposited mass from shedders, decay each epoch, and are picked up by later agents with a pickup probability and transfer
fraction. Quarantined agents in cabin corridors skip public fomite pickup, while cabin-mate exposure remains possible. Food contamination is a separate pool for dining-type zones. A small
fraction of shedding deposits into food pools and a configured fraction of that pool is ingested by agents in enabled food zones. Environmental-source transmission supports pathogens such as
Legionella by delivering dose from a zone-scoped or global environmental reservoir even in the absence of infected human shedders.

\section{CONTAM HVAC coupling}
The CONTAM branch is implemented by \texttt{engines/py\_contam\_bridge.py} as a native Python bridge to multizone airflow and contaminant transport. It reads \texttt{air\_flow\_paths.json},
builds zone nodes from the platform spatial layout, and constructs airflow paths from three sources: HVAC recirculation, cross-zone links, and passive adjacency. The star-topology HVAC
representation sends room return air to a virtual AHU plenum and then supplies filtered recirculated air back to rooms. Outdoor-air fraction removes a fraction of return pathogen mass before
supply, and filter efficiency applies to ducted supply paths.

For zone $z$, the bridge solves a well-mixed contaminant ODE analytically at each epoch:
\begin{equation}
\frac{dM_{z,p}}{dt}=S_{z,p}(t)-k_z M_{z,p}(t),
\end{equation}
\begin{equation}
M_{z,p}(t+\Delta t)=M_{z,p}(t)e^{-k_z\Delta t}+\frac{S_{z,p}(t)}{k_z}
\left(1-e^{-k_z\Delta t}\right),
\label{eq:contam-ode}
\end{equation}
where $S_{z,p}$ is the filtered inflow source rate and
\begin{equation}
k_z=\sum_{z'}\frac{Q_{z\rightarrow z'}}{V_z}+\lambda_p.
\end{equation}
The terms $Q_{z\rightarrow z'}$, $V_z$, and $\lambda_p$ are volumetric flow, zone volume, and natural decay or settling rate. This branch is used when airborne route attribution or
native-versus-CONTAM validation is part of the experiment. For campaign-scale runs where HVAC is not the tested mechanism, the native transport branch preserves tractable runtime.

\section{Observation engine and diagnostic cascade}
CTB separates infection from observation. The observation engine samples the simulated ship through instruments rather than assuming that all infections are immediately known. The instrument
set includes continuous air sampling, targeted surface swabs, wastewater sequencing grids, clinical rapid diagnostics, clinical qPCR, clinical microbiology, long-read verification, wearable
physiological monitoring, and fleet-level wearable summaries. Each instrument produces raw observations such as Ct values, detection booleans, anomaly scores, clinical positives, or
microbiome disruption values. Those raw observations are then converted to stoplight states.

A generic observation model for modality $m$ can be written as
\begin{equation}
Y_m(t)=g_m\{\mathcal{S}(t-\delta_m),\theta_m\}+\epsilon_m(t),
\end{equation}
where $\delta_m$ is reporting delay, $\theta_m$ contains sensitivity, specificity, coverage, and threshold parameters, and $\epsilon_m$ is instrument noise. The stoplight function is
\begin{equation}
Z_m(t)=h_m(Y_m(t))\in\{\mathrm{GREEN},\mathrm{AMBER},\mathrm{RED}\}.
\end{equation}
For Ct-based instruments, low Ct values map to stronger stoplights when detection is positive. For wastewater and microbiome disruption, anomaly scores map to stoplights. For clinical RDTs,
positives map directly to RED. For syndromic surveillance, the sick-call count is compared with configured AMBER and RED thresholds.

The diagnostic cascade uses those stoplights to represent escalation across information sources. A virtual \texttt{detection\_escalation} instrument aggregates syndromic, wearable,
environmental, and clinical modes. For a set of modes $\mathcal{M}$, the integrated level can be represented as
\begin{equation}
L(t)=\max_{m\in\mathcal{M}} Z_m(t),
\end{equation}
with the ordering GREEN $<$ AMBER $<$ RED. Protocols can require multiple affected modes, zones, or agents before firing. This design permits onboard response to depend on concordant signals
rather than on a single perfect observer.

\section{Reactive protocol engine}
The reactive protocol engine is implemented in \texttt{crusher\_labs/protocol\_engine.py} and configured through protocol JSON files validated by schema contracts. At the end of each epoch,
it recomputes stoplights from raw results, checks SOP trigger conditions, applies escalation gates, observes activation delays, activates or deactivates SOPs, merges modifiers, and debits
costs through the cost ledger. This makes response a function of evidence generated by the simulated ship. The model rule is no hardcoded epoch schedules for outbreak response; interventions
fire because instrument or attack-rate conditions are met.

A protocol $s$ is eligible when its trigger condition is satisfied:
\begin{equation}
\mathbb{1}_{s}^{\mathrm{trigger}}(t)=
\mathbf{1}\left\{\sum_{u\in U_s}\mathbf{1}\left[Z_{m_s,u}(t)\ge z_s\right]\ge n_s\right\},
\end{equation}
where $m_s$ is the instrument class, $U_s$ is the set of zones, agents, or modes evaluated by the trigger, $z_s$ is the required stoplight level, and $n_s$ is the minimum affected count. If
the SOP has an escalation gate, ship status must also exceed the required rank. If it has an activation delay $\ell_s$, modifiers apply only after the trigger has persisted long enough:
\begin{equation}
A_s(t)=\mathbf{1}\{\mathbb{1}_{s}^{\mathrm{trigger}}(t-\ell_s)=1,\;E(t)\ge E_s^{\min}\}.
\end{equation}

SOP modifiers are ordinary model parameters exposed through the protocol layer. Examples include \texttt{direct\_contact\_scalar},\\ \texttt{droplet\_scalar},
\texttt{hvac\_filter\_efficiency\_override}, \texttt{close\_zones}, diagnostic cadence changes, and PPE or cleaning effects. Scalar reductions use the most restrictive value when several
protocols are active; list-based closures are unioned. The cost ledger records both one-time activation costs and recurring per-epoch costs. This cost accounting is important because
surveillance is evaluated not only by attack-rate reduction but also by the decision quality of when expensive escalation occurs.

VSP-style escalation is represented as a threshold rule on observed illness rather than on true infection. Let $\mathrm{AR}_{\mathrm{cum}}(t)$ be cumulative reportable acute-gastroenteritis
attack rate. The VSP trigger is
\begin{equation}
\mathrm{escalate}(t)=\mathbf{1}\{\mathrm{AR}_{\mathrm{cum}}(t)>\theta_{\mathrm{VSP}}\}.
\label{eq:vsp-trigger}
\end{equation}
The baseline calibration uses the VSP reporting threshold as a response backstop. Because the trigger is tied to observed illness, the same biological parameterization can produce a high
uncontrolled attack rate and a much lower public record after response.

\section{Overdispersion and superspreading}
Superspreading is represented as an emergent distribution of secondary cases rather than as a fixed label attached to a subset of agents. The conventional summary is the negative-binomial
offspring distribution
\begin{equation}
Z_i\sim\operatorname{NegBin}(R,k),\qquad
\operatorname{Var}(Z_i)=R+\frac{R^2}{k},
\end{equation}
where small $k$ indicates strong overdispersion \citep{LloydSmith2005,Endo2020Overdispersion}. For SARS-CoV-2, early cluster analyses estimated strong overdispersion, and reviews emphasize
that large events arise from the concurrence of infectiousness, behavior, environment, and opportunity \citep{Althouse2020Superspreading}. At the same time, a fitted negative-binomial $k$
can conflate network structure, observation bias, temporal infectiousness, intervention timing, and setting mixtures \citep{Zhao2022NegativeBinomial}.

ABMs capture overdispersion naturally when their contact and environment layers are sufficiently heterogeneous. In CTB, large offspring counts can arise when an infectious agent has high
shedding near peak infectiousness, works as crew in a dining or galley zone, attends a crowded venue, shares a cabin, contaminates a food or surface pool, or occupies a poorly ventilated
zone that feeds downstream airflow. The heterogeneous-zone-dose mode adds residual within-zone variability while preserving the mean dose. The preferred calibration workflow is therefore
mechanism first and offspring-distribution check second. If the simulated zero-transmitter fraction, upper-tail secondary-case counts, largest event size, and extinction probability already
match data, adding an independent superspreader multiplier would double count heterogeneity.

This point is especially relevant for cruise ships. A deterministic or homogeneous compartmental model can reproduce a mean epidemic curve while missing fade-outs and rare large outbreaks. A
stratified model can separate passengers, crew, and decks, but it still averages within strata. CTB records exposure matrices and transmission events by pathway, zone, and pathogen, allowing
the model to report distributions of secondary cases and route attributions rather than only final attack rates. Those outputs are the appropriate validation targets when line lists, cabin
maps, venue records, or sequence data become available.

\section{Multi-pathogen support}
Multi-pathogen support is implemented through pathogen profiles, route weights, dose-response parameters, environmental settings, food-contamination settings, and pathogen-specific infection
state. The schema layer defines contracts for profiles, protocol files, run specifications, instrument parameters, decision actions, and voyage configurations. These contracts are part of
the architecture: they prevent a campaign from silently changing the meaning of a field across pathogens or platforms.

For pathogen $p$, the profile defines a parameter set
\begin{equation}
\Theta_p=\{\alpha_p,\beta_p,\kappa_p,\doseadj_p,w_{p,r},\lambda_p,\phi_{p,r},\chi_p\},
\end{equation}
where $\alpha_p$ and $\beta_p$ parameterize beta-Poisson dose response, $\kappa_p$ parameterizes exponential response when used, $w_{p,r}$ are route weights, $\lambda_p$ is environmental
decay, $\phi_{p,r}$ are route-specific deposition or delivery fractions, and $\chi_p$ contains pathogen-specific clinical and environmental options. Norovirus uses contact, fomite, food, and
wastewater-relevant observation. SARS-CoV-2 and measles emphasize droplet and airborne pathways, diagnostic timing, and respiratory presentation. Legionella emphasizes environmental-source
and HVAC-linked exposure rather than person-to-person spread.

The monograph campaigns use this architecture across 4 pathogens and 4 vessel platforms. The same epoch loop, observation engine, protocol engine, and telemetry contracts are reused, but the
biological parameters and route weights change. This is what makes CTB a platform for comparing surveillance and response logic rather than a single-pathogen curve-fitting exercise. The
architecture also makes explicit where uncertainty remains: route-specific dose conversion, environmental signal calibration, individual compliance, and the mapping from simulated pathogen
mass to measured Ct or sequencing counts. Those uncertainties are carried into later chapters as calibration and validation problems rather than hidden inside a single effective transmission
rate.
